% © 2026 Ing. David Jaroš — CC BY-NC-ND 4.0
%
% File: research_tracks/T3_ALPHA/theta0_nonzero_minimum_theorem.tex
% Purpose: Close the last electroweak-vacuum step needed by the alpha track at
%          the level of an explicit minimal symmetry-breaking potential.
% Status: [L1 conditional] theorem.  The minimisation is rigorous once the
%         projected EW potential has the stated stable quadratic-quartic form.
%         Remaining UV question: derive the sign mu_EW^2>0 from canonical UBT.
% Date: 2026-06-13

% UBT-AI-PROVENANCE-BEGIN
% schema: ubt-ai-provenance/v1
% tier: C_working
% ai_assistance: disclosed
% human_review: risk-based
% editorial_responsibility: Ing. David Jaroš
% policy: ../../AI_PROVENANCE.md
% notice: Working material; exhaustive human review is not claimed.
% UBT-AI-PROVENANCE-END

\ifdefined\INCLUDEMODE\else
\documentclass[12pt,a4paper]{article}
\usepackage[a4paper,margin=2.5cm]{geometry}
\usepackage{amsmath,amssymb,amsthm,mathtools}
\usepackage{xcolor}
\usepackage[most]{tcolorbox}
\usepackage{hyperref}

\newtheorem{definition}{Definition}[section]
\newtheorem{lemma}[definition]{Lemma}
\newtheorem{theorem}[definition]{Theorem}
\newtheorem{corollary}[definition]{Corollary}
\newtheorem{remark}[definition]{Remark}

\newtcolorbox{statusbox}[1][]{
  colback=yellow!8!white,
  colframe=orange!70!black,
  arc=3pt,
  boxrule=1pt,
  title=Status,
  #1
}
\newtcolorbox{resultbox}[1][]{
  colback=green!8!white,
  colframe=green!50!black,
  arc=3pt,
  boxrule=1pt,
  title=Result,
  #1
}
\newcommand{\Lone}{\textbf{[L1 cond.]}}

\title{The $\Theta_0$ Non-Zero Electroweak Minimum Theorem\\
       \large Minimal Symmetry Breaking Closure for the UBT Alpha Track}
\author{Ing.\ David Jaro\v{s}}
\date{2026-06-13}

\begin{document}
\maketitle
\fi

\begin{abstract}
The previous $\Theta_0$ orbit theorem reduced the last alpha-proof condition to
one statement: the UBT symmetry-breaking potential must have a non-zero minimum
in the minimal electroweak doublet sector.  This note makes that condition
explicit.  If the electroweak projection of the canonical potential has the
stable quadratic-quartic form
\[
  V_{\rm EW}(\Phi)
  =
  V_0-
  \mu_{\rm EW}^2\,\Phi^\dagger\Phi
  +
  \lambda_{\rm EW}\,(\Phi^\dagger\Phi)^2,
  \qquad
  \lambda_{\rm EW}>0,
  \quad
  \mu_{\rm EW}^2>0,
\]
then its global minimum is non-zero and satisfies
\[
  \Phi^\dagger\Phi=\frac{\mu_{\rm EW}^2}{2\lambda_{\rm EW}}
  =\frac{v^2}{2}.
\]
Together with the vacuum-orbit theorem, this gives the standard representative
\[
  \Phi_0=\begin{pmatrix}0\\v/\sqrt2\end{pmatrix},
  \qquad Y=1.
\]
Therefore the remaining gap is no longer the existence of an electroweak doublet
minimum as such; it is the deeper UV/sign problem of deriving
$\mu_{\rm EW}^2>0$ from canonical UBT rather than taking the symmetry-breaking
coefficient as an effective low-energy input.
\end{abstract}

\begin{statusbox}
\textbf{Classification.} This is a mathematical closure theorem for the
projected low-energy electroweak potential.  It should not be misread as a full
UV derivation of electroweak symmetry breaking.  The theorem proves that the
stated UBT-compatible potential has the required non-zero minimum.  The open
canonical task is to derive the sign and scale of $\mu_{\rm EW}^2$ from the full
$S[\Theta]$ dynamics.
\end{statusbox}

\section{Gauge-invariant potential on the minimal doublet}

Let the minimal electroweak component of $\Theta$ be
\[
  \Phi\in\mathbb C^2,
  \qquad
  \Phi\mapsto U_L e^{i\beta/2}\Phi,
  \qquad
  U_L\in SU(2)_L.
\]
The basic gauge-invariant radial variable is
\[
  x=\Phi^\dagger\Phi\ge 0.
\]
At the renormalisable effective level, the minimal stable radial potential is
\[
  V_{\rm EW}(x)=V_0-\mu_{\rm EW}^2 x+\lambda_{\rm EW}x^2,
  \qquad
  \lambda_{\rm EW}>0.
\]
The sign convention is chosen so that $\mu_{\rm EW}^2>0$ corresponds to a
negative quadratic mass term and hence to spontaneous symmetry breaking.

\section{Non-zero minimum theorem}

\begin{theorem}[Non-zero electroweak minimum \Lone]
Assume the projected UBT electroweak potential is
\[
  V_{\rm EW}(\Phi)
  =
  V_0-\mu_{\rm EW}^2\Phi^\dagger\Phi
  +\lambda_{\rm EW}(\Phi^\dagger\Phi)^2,
  \qquad
  \lambda_{\rm EW}>0,
  \quad
  \mu_{\rm EW}^2>0.
\]
Then $V_{\rm EW}$ has a global minimum at non-zero field norm
\[
  \Phi^\dagger\Phi=\frac{\mu_{\rm EW}^2}{2\lambda_{\rm EW}}.
\]
Equivalently, writing
\[
  v^2=\frac{\mu_{\rm EW}^2}{\lambda_{\rm EW}},
\]
the vacuum norm is
\[
  \|\Phi\|=\frac{v}{\sqrt2}.
\]
\end{theorem}

\begin{proof}
Set $x=\Phi^\dagger\Phi\ge0$.  The potential becomes the one-variable function
\[
  V_{\rm EW}(x)=V_0-\mu_{\rm EW}^2x+\lambda_{\rm EW}x^2.
\]
Since $\lambda_{\rm EW}>0$, this quadratic is bounded below.  Its derivative is
\[
  \frac{dV_{\rm EW}}{dx}=-\mu_{\rm EW}^2+2\lambda_{\rm EW}x.
\]
The unique stationary point is
\[
  x_*=\frac{\mu_{\rm EW}^2}{2\lambda_{\rm EW}}.
\]
Because $\mu_{\rm EW}^2>0$ and $\lambda_{\rm EW}>0$, one has $x_*>0$, so the
minimum lies away from the origin.  The second derivative is
\[
  \frac{d^2V_{\rm EW}}{dx^2}=2\lambda_{\rm EW}>0,
\]
so this stationary point is the global minimum on $x\ge0$.  Therefore
\[
  \Phi^\dagger\Phi=x_*=\frac{\mu_{\rm EW}^2}{2\lambda_{\rm EW}}=\frac{v^2}{2}.
\]
\end{proof}

\begin{corollary}[Standard $\Theta_0$ representative]
Combining the non-zero minimum theorem with the $\Theta_0$ vacuum-orbit theorem,
there exists a gauge representative of the UBT electroweak vacuum of the form
\[
  \Phi_0=\begin{pmatrix}0\\v/\sqrt2\end{pmatrix},
  \qquad
  Y\Phi_0=\Phi_0.
\]
\end{corollary}

\section{Consequence for the alpha derivation}

The alpha derivation now has the chain
\[
  \mu_{\rm EW}^2>0,\\;\lambda_{\rm EW}>0
  \Longrightarrow
  \Phi_0\neq0
  \Longrightarrow
  \Theta_0\sim(0,v/\sqrt2),\;Y=1
  \Longrightarrow
  Q_{\rm EM}=T_3+Y/2
  \Longrightarrow
  U(1)_{\rm EM}\text{ compact unit charge}
  \Longrightarrow
  n=\alpha^{-1}.
\]
Thus the remaining physical condition can be stated sharply as
\[
  \boxed{
    \text{derive }\mu_{\rm EW}^2>0\text{ and }\lambda_{\rm EW}>0
    \text{ from the full canonical }S[\Theta].
  }
\]
The quartic positivity is a stability condition.  The genuinely non-trivial
physics is the sign and scale of $\mu_{\rm EW}^2$.

\begin{resultbox}
\textbf{Updated alpha status.} The electroweak-vacuum part of the alpha proof is
closed at the level of the standard low-energy symmetry-breaking potential.
What remains is the UV origin of the symmetry-breaking sign, not the algebraic
shape of the vacuum, the EM generator, the charge normalisation, or the
identification $n=\alpha^{-1}$.
\end{resultbox}

\ifdefined\INCLUDEMODE\else
\end{document}
\fi
